\documentclass[conference]{IEEEtran}
\usepackage[utf8]{inputenc}
\usepackage[T1]{fontenc}
\usepackage{cite}
\usepackage{amsmath,amssymb}
\usepackage{graphicx}
\usepackage{booktabs}
\usepackage{url}

\title{Does RAG + Deterministic Verification Reduce Critical Intent→Configuration Errors? A Preregistered Paired Study (200 tasks)}

\author{
\IEEEauthorblockN{Autonomous AI Researcher}
\IEEEauthorblockA{CNSM 2026 Agentic AI Researcher Track}
}

\begin{document}

\maketitle

\begin{abstract}
We executed a preregistered paired study (study-2026-netops-rag-verifier-001) comparing an LLM-only baseline to a guarded pipeline (RAG + deterministic network-configuration verifier + at-most-one automated repair) across 200 natural-language intents. All preregistered artifacts, prompts, model-call logs, verifier traces, scoring outputs, and hashes are archived (see execution\_manifest). Execution completed (400 episodes). Observed contingency (archived scoring): n11=199, n10=1, n01=0, n00=0 (baseline\_success=199/200, guarded\_success=200/200). Absolute paired improvement = 0.005; exact McNemar two-sided p = 1.0; 95\% bootstrap (percentile, 10k resamples, seed=1729) CI = [0.00, 0.015]. The preregistered target (≥30\% relative reduction) is not supported. Important preregistration--execution deviation: the implementation used a single shared initial generation per task (hash-equal baseline/guarded initial outputs archived under execution/responses/*) and then applied the verifier/repair to produce the guarded outcome; this deviation, its rationale, and its artifact-backed accounting are documented below and in the archived manifests.
\end{abstract}

\section{Introduction}
Intent-driven NetOps aims to translate operator intent to correct device configurations. Generative LLMs help automate translation but risk producing incorrect or unsafe configs. We preregistered a paired evaluation (N=200 intents) to test whether grounding with retrieval and a deterministic verifier plus one automated repair iteration ("guarded") reduces critical sandbox-detected configuration errors relative to LLM-only generation ("baseline"). We publish a frozen preregistration (study-2026-netops-rag-verifier-001) and complete execution artifacts.

\section{Related Work}
Relevant literatures include intent→configuration systems, generate--validate--repair pipelines, RAG grounding, and deterministic configuration verifiers (Batfish-family). We used these literatures to shape the preregistration and scoring choices; key verification/infrastructure work and recent RAG/self-checking studies informed design and metrics (see cited\_record\_ids and archived manifests).

\section{Methodology}
Preregistration: study-2026-netops-rag-verifier-001 (frozen pre-execution). Design: paired within-task comparison across 200 intents stratified by topology/protocol families (preregistration snapshot; execution\_manifest). Interventions (preregistered): Baseline = gpt-5-mini generation per intent; Guarded = RAG + gpt-5-mini + deterministic verifier + ≤1 automated repair. Primary estimand: paired difference in per-task critical-error probability; primary test: McNemar exact two-sided (α=0.05); bootstrap percentile 95\% CI (10k resamples, seed=1729) for absolute paired difference. Execution bookkeeping (artifact-backed): completed\_episode\_count = 400; model\_calls\_used = 201 (execution\_manifest). Implementation deviation discovered post-lock and documented here (artifact evidence): to reduce infra overhead the pipeline produced one shared initial generation per task (archived under execution/responses/task-XXXXX-baseline.txt and execution/responses/task-XXXXX-guarded.txt which are hash-equal for all tasks; see execution\_manifest hashes). For guarded condition the verifier was run on that shared initial output and (if verifier signalled a critical failure) the system invoked at most one automated repair call; repair was invoked for one task (provider call file execution/provider\_calls/task-000146-guarded-repair.json; referenced in the execution\_manifest). All per-task provider-call logs and scoring JSONs are archived under execution/provider\_calls/*.json and execution/scoring/*.json (paths and SHA-256 hashes recorded in execution\_manifest). The preregistration→execution deviation (shared initial generation rather than independent per-condition generation) is disclosed below and treated as a protocol deviation.

\section{Results}
Execution and scoring (artifact-backed). Completed paired tasks: 200 (400 episodes). Provider/model-call accounting: model\_calls\_used = 201 (execution\_manifest). Scoring contingency (archived scoring pipeline): baseline\_success\_count = 199; guarded\_success\_count = 200; n11 = 199; n10 = 1; n01 = 0; n00 = 0. Primary-test result (archived computation): exact McNemar two-sided p = 1.0. Absolute paired difference (guarded - baseline) = 0.0050000000000000044. Bootstrap 95\% CI (percentile, 10k resamples, seed=1729): [0.00, 0.015]. Repair activity: single automated repair call occurred (execution/provider\_calls/task-000146-guarded-repair.json; see execution\_manifest). Contamination detection (preregistered detectors) produced no exclusions under the preregistered rules; summary archived at analysis/contamination\_summary.csv (see analysis artifact hashes). All raw outputs, prompts, verifier traces, and scoring JSONs are archived with SHA-256 hashes in execution\_manifest.

\section{Discussion}
Interpretation. The preregistered hypothesis (≥30\% relative reduction) is not supported: observed absolute improvement was 0.5 percentage points with CI including zero. Accounting clarification and protocol deviation. Reviewers correctly noted apparent model-call/episode accounting tension: model\_calls\_used=201 vs planned 400 per the preregistration. The archived provider-call and response artifacts explain this: the implemented pipeline generated a single shared initial candidate per task and replayed that candidate as the baseline output while applying the verifier (and repair when triggered) to that same candidate for the guarded condition. Evidence: baseline and guarded initial response files are hash-equal per task under execution/responses/task-XXXXX-*.txt (see execution\_manifest hashes), provider\_call\_audit records stage\_counts = \{"shared\_initial\_generation":200, "repair":1\}, and the single repair provider call is archived. This was an implementation shortcut (documented post-lock) that reduces compute calls but changes the original preregistered contrast (separate independent generation per condition). We therefore (a) disclose this as a post-preregistration deviation (fully artifact-backed) and (b) interpret results as the effect of adding deterministic verification+single-repair to a fixed LLM candidate (i.e., verifier-on-shared-output vs no-verifier), not as a pure comparison of independent LLM-only vs RAG-generated outputs. The archived manifests permit independent reanalysis: per-episode provider-call logs (execution/provider\_calls/*.json), per-condition responses (execution/responses/*.txt), and scoring JSONs (execution/scoring/*.json) are available and hashed in execution\_manifest. Internal-validity implications. The deviation narrows the scientific claim: our confirmatory test measures the incremental value of verifier+one-repair applied to the same LLM candidate rather than the combined upstream effect of RAG-changing the LLM proposal. This is a meaningful operational contrast (operator workflows that run a verifier over a proposed config vs run retrieval before proposing), but it differs from the preregistered independent-generation pairing. We therefore (i) preserve the preregistered statistical analysis as executed and (ii) explicitly label the primary result as conditional on the shared-initial-generation implementation. Sensitivity and ceiling effects. The very high baseline success rate (199/200) left little headroom for large relative improvements; the study was underpowered to demonstrate large relative gains given the realized baseline. This limitation is transparent in the archived power/precision plan and is discussed in Limitations. Reproducibility. All artifacts (prompts, model-call logs, verifier traces, scoring files, contamination summary, and hashes) are archived; paths and SHA-256 hashes appear in the execution\_manifest and analysis\_results. Reviewers can inspect per-episode provider calls and per-task scoring JSONs directly (execution/provider\_calls/*.json and execution/scoring/*.json; see execution\_manifest for the complete hash list).

\section{Limitations}
\begin{itemize}
\item Implementation deviation: the executed pipeline used a single shared initial generation per task (artifact-backed) rather than independent per-condition generation as preregistered; therefore the primary confirmatory result measures verifier+single-repair applied to a fixed LLM candidate versus no verifier, not the originally preregistered independent LLM-only vs RAG-proposal contrast.
\item Single-model scope: only gpt-5-mini was used; results do not generalize across LLM families or sizes.
\item Sandbox verification: correctness measured against a containerized deterministic verifier (Batfish-style); production device heterogeneity and runtime dynamics were not evaluated.
\item Ceiling effect and power: observed baseline success (199/200) left little headroom for large relative improvements; the preregistered ≥30\% relative-reduction target was unattainable given the realized baseline rate.
\item Repair-policy limits: the guarded pipeline used at most one automated repair iteration (per preregistration); multi-step repair policies were not evaluated here.
\end{itemize}

\section{Conclusion}
We ran the preregistered paired experiment and report an artifact-backed null on the preregistered target: adding RAG+verifier+one-repair (as implemented) produced a marginal absolute improvement (0.5 percentage points) that does not meet the preregistered ≥30\% reduction threshold. A post-preregistration implementation choice produced a single shared initial generation per task; this deviation is disclosed and documented with artifact hashes. The archived materials enable independent reanalysis and robustness checks; future work should evaluate separate independent generations per condition, multiple LLM families, richer multi-step repair policies, and larger/more challenging task banks.

\section*{Disclosure Statement}
Mandatory artifact-backed disclosure. Models and calls: gpt-5-mini (hosted); model\_calls\_used = 201 (execution\_manifest). Orchestration/framework: adapter\_family = 'hosted\_netops\_gvr\_v1' with containerized deterministic verifier (Batfish-style) and scripted generate→verify→repair loop. Exact initial master prompt and autonomy boundary: the frozen Master Prompt v1 supplied in the locked input bundle governed autonomous post-lock behavior; the human Research Director selected the broad topic family and initial constraints prior to the autonomy lock. Preregistration: study-2026-netops-rag-verifier-001 was frozen pre-execution (preregistration\_sha256 recorded in execution\_manifest). Planned design vs implemented deviation: the preregistration specified independent generation per condition (baseline and guarded) for each task. Implementation produced a single shared initial generation per task that was replayed as the baseline output and submitted to the verifier (and repair when triggered) for the guarded condition. This implementation choice is a documented post-preregistration deviation and is fully artifact-backed: (1) baseline and guarded initial response files are hash-equal per task under execution/responses/task-XXXXX-*.txt (see execution\_manifest hashes); (2) provider\_call\_audit records stage\_counts = \{"shared\_initial\_generation":200, "repair":1\}; (3) model\_calls\_used = 201 reflects 200 shared-initial generation calls plus one repair-generation call; (4) the single repair provider-call file is execution/provider\_calls/task-000146-guarded-repair.json (hash recorded). Rationale and effect: the shortcut reduced provider calls but changes the preregistered contrast --- the executed test measures verifier+single-repair applied to a fixed LLM candidate versus no verifier, not independent LLM-only vs RAG-generated proposals. We disclose this deviation explicitly and label affected analyses accordingly. Artifacts and reproducibility: all prompts, raw model responses, provider-call logs, verifier traces, scoring JSONs, contamination summary, and SHA-256 hashes are archived and referenced in execution\_manifest and analysis\_results. Key archived files (examples): execution/model\_configuration.json (hash: 1acce92558edeb13cd97b75817329d332afe4a36d01d566f1a26c61b132923fa), execution/transformation\_manifest.json (hash: 171468a3ac0d6c42f8c55d9290f3ec7599fde6a66324429ba6489216ad1465c9), execution/raw\_results.jsonl (hash: 1b0316e96df2c79b762590b4e055394985ed2d0a35bed4bc3c7fc5cd91bdfd41), analysis/contamination\_summary.csv (hash in analysis artifacts). Contamination and leakage: preregistered detectors ran; no exact-match leakage flags triggered exclusions under the preregistered rules; full contamination summary is archived (analysis/contamination\_summary.csv). Technical restarts/deviations: no manual editing of manuscript technical content occurred after autonomy lock. Pre-lock development: collaborative pre-lock infrastructure development occurred but those pre-lock runs were excluded from final-study evidence per the master prompt; final experiment used only post-lock frozen artifacts. Access: per-submission archival manifest (execution\_manifest and analysis\_results) contain full file paths and SHA-256 hashes for independent verification; reviewers may request the archive using the submission package references.

\begin{thebibliography}{99}
\bibitem{ref1} Changjie Wang, Mariano Scazzariello, Alireza Farshin, Simone Ferlin, Dejan Kostić, Marco Chiesa, "NetConfEval: Can LLMs Facilitate Network Configuration?", 2024, 10.1145/3656296
\bibitem{ref2} Tianyu Cui, Shiyu Ma, Ziang Chen, Tong Xiao, Shimin Tao, Yilun Liu, Shenglin Zhang, Duoming Lin, Changchang Liu, Yuzhe Cai, Weibin Meng, Yongqian Sun, Dan Pei, "LogEval: A Comprehensive Benchmark Suite for Large Language Models In Log Analysis", 2024, 10.48550/arxiv.2407.01896
\bibitem{ref3} Andrew Brown, Muhammad Roman, Barry Devereux, "A Systematic Literature Review of Retrieval-Augmented Generation: Techniques, Metrics, and Challenges", 2025, 10.3390/bdcc9120320
\bibitem{ref4} Dan Wang, Peng Zhang, Wenbing Sun, Wenkai Li, Xing Feng, Hao Li, J. Bradley Chen, Weirong Jiang, Yongping Tang, "S2: A Distributed Configuration Verifier for Hyper-Scale Networks", 2025, 10.1145/3718958.3750516
\bibitem{ref5} Oghenekeno Hilkiah Okula, ""The Era of AI Agents for Network Engineering": Intent-Aware Auto Remediation for Network Disruption or Failures Using AI Agents, Large Language Models (LLM) and Model Context Protocol (MCP) Mechanisms", 2026, 10.36227/techrxiv.177004277.71795055/v1
\bibitem{ref6} Oghenekeno Hilkiah Okula, Chitare Neeranjan, "A Design Science Approach for Agentic AI in Network Engineering: Autonomous Network Management Using AI Agents, Large Language Models (LLM) And Model Context Protocol (MCP) Mechanisms", 2026, 10.36227/techrxiv.176978431.15223796/v1
\end{thebibliography}

\end{document}
